\documentclass[11pt, a4paper]{article}
\usepackage{amsmath, amssymb, booktabs, geometry, hyperref, xcolor, listings}
\geometry{margin=2.5cm}
\hypersetup{colorlinks=true, linkcolor=blue}
\lstset{basicstyle=\ttfamily\small, backgroundcolor=\color{gray!10}, frame=single}

\title{\textbf{Multivariate Monte Carlo Simulation}\\[4pt]
\large Portfolio VaR Engine --- Model Documentation}
\author{Risk Modelling --- Trade Republic}
\date{April 2026}

\begin{document}
\maketitle

%-----------------------------------------------------------------------
\section{Overview}

The Monte Carlo (MC) engine generates a large number of independent
synthetic return paths for all assets in the portfolio simultaneously,
using a multivariate distribution calibrated to historical data.
VaR and CVaR are then estimated empirically from the simulated portfolio
P\&L distribution. This is the only engine that natively captures
cross-asset dependence via a full covariance matrix.

\textbf{Engine key:} \texttt{'mc'} \\
\textbf{Sources:}
\begin{itemize}
  \item \texttt{risk\_models/mc\_models.py}, lines 16--96 (\texttt{mc\_simulate})
  \item \texttt{risk\_analytics/risk\_engines.py}, lines 154--186 (VaR/CVaR aggregation)
  \item \texttt{risk\_models/support\_models.py}, lines 49--61 (\texttt{ewma\_covariance})
\end{itemize}

%-----------------------------------------------------------------------
\section{Covariance Matrix Estimation}

Three options are available for the input covariance matrix
$\hat{\Sigma} \in \mathbb{R}^{N \times N}$ (where $N$ is the number of assets):

\subsection{EWMA Covariance (default: \texttt{ewma\_cov=True})}

Following RiskMetrics, the EWMA covariance matrix over a window of $T$
observations is:
\begin{equation}
\hat{\Sigma}_{\text{EWMA}}
= \frac{\sum_{k=0}^{T-1} \lambda^k \, (\mathbf{r}_{T-k} - \bar{\mathbf{r}})
                                      (\mathbf{r}_{T-k} - \bar{\mathbf{r}})^\top}
       {\sum_{k=0}^{T-1} \lambda^k}
\label{eq:ewma_cov}
\end{equation}
where $\bar{\mathbf{r}}$ is the sample mean vector and $\lambda = 0.94$
is the decay factor. More recent observations receive exponentially higher
weight.

\subsection{Sample (Historical) Covariance (\texttt{ewma\_cov=False})}

\begin{equation}
\hat{\Sigma}_{\text{hist}} = \frac{1}{T-1}
  \sum_{t=1}^{T} (\mathbf{r}_t - \bar{\mathbf{r}})(\mathbf{r}_t - \bar{\mathbf{r}})^\top
\end{equation}

\subsection{Maximum Likelihood Estimator (\texttt{mlest\_cov=True})}

Uses \texttt{sklearn.covariance.EmpiricalCovariance} on 500 Monte Carlo
samples drawn from a multivariate Normal with the historical covariance.
This option is available for stress scenarios where \texttt{stress\_means}
are specified.

%-----------------------------------------------------------------------
\section{Simulation Procedure}

\subsection{Drawing Simulated Paths}

For each of $n$ simulation runs $i = 1, \ldots, n$, draw a matrix of
$T_{\text{sim}}$ multivariate daily returns:

\textbf{Multivariate Normal:}
\begin{equation}
\mathbf{r}^{(i)} \overset{i.i.d.}{\sim}
\mathcal{N}_N\!\left(\boldsymbol{\mu},\, \hat{\Sigma}\right),
\quad \mathbf{r}^{(i)} \in \mathbb{R}^{T_{\text{sim}} \times N}
\end{equation}

\textbf{Multivariate Student-$t$:}
\begin{equation}
\mathbf{r}^{(i)} \overset{i.i.d.}{\sim}
t_N\!\left(\boldsymbol{\mu},\, \hat{\Sigma},\, \nu\right),
\quad \mathbf{r}^{(i)} \in \mathbb{R}^{T_{\text{sim}} \times N}
\end{equation}
where $\boldsymbol{\mu}$ is the historical mean vector and $\nu$ is the
degrees-of-freedom parameter (default $\nu = 5$).

The Student-$t$ option produces heavier-tailed joint distributions and
is therefore more conservative for tail risk estimation.

\subsection{Portfolio Return Paths}

For each simulation $i$, the portfolio return path is:
\begin{equation}
r_{pf,t}^{(i)} = \mathbf{w}^\top \mathbf{r}_t^{(i)}, \quad t = 1,\ldots, T_{\text{sim}}
\end{equation}
giving a simulated P\&L vector $\mathbf{r}_{pf}^{(i)} \in \mathbb{R}^{T_{\text{sim}}}$
for each simulation run.

%-----------------------------------------------------------------------
\section{VaR and CVaR Estimation}

\subsection{Value-at-Risk}

For each simulation run $i$, compute the empirical $\alpha$-quantile of
the portfolio return path:
\begin{equation}
\text{VaR}_\alpha^{(i)} = -Q_\alpha\!\left(\mathbf{r}_{pf}^{(i)}\right)
\end{equation}
The reported VaR is the \textbf{median} across all $n$ simulation runs:
\begin{equation}
\widehat{\text{VaR}}_\alpha
= \operatorname{median}_{i=1,\ldots,n}\!\left\{\text{VaR}_\alpha^{(i)}\right\}
\label{eq:mc_var}
\end{equation}
Using the median (rather than the mean) provides robustness against
extreme outlier simulations.

\subsection{Conditional Value-at-Risk (Expected Shortfall)}

For each simulation run $i$, CVaR is the mean of returns that fall at or
below the $\alpha$-quantile:
\begin{equation}
\text{CVaR}_\alpha^{(i)}
= -\frac{1}{|\mathcal{S}_\alpha^{(i)}|}
  \sum_{t \in \mathcal{S}_\alpha^{(i)}} r_{pf,t}^{(i)},
\qquad
\mathcal{S}_\alpha^{(i)} = \left\{t : r_{pf,t}^{(i)} \leq Q_\alpha\!\left(\mathbf{r}_{pf}^{(i)}\right)\right\}
\end{equation}
The reported CVaR is again the median across simulations:
\begin{equation}
\widehat{\text{CVaR}}_\alpha
= \operatorname{median}_{i=1,\ldots,n}\!\left\{\text{CVaR}_\alpha^{(i)}\right\}
\label{eq:mc_cvar}
\end{equation}

\subsection{Holding-Period Scaling}

\begin{equation}
\widehat{\text{VaR}}_\alpha^{(H)} = \widehat{\text{VaR}}_\alpha^{(1)} \cdot \sqrt{H},
\qquad
\widehat{\text{CVaR}}_\alpha^{(H)} = \widehat{\text{CVaR}}_\alpha^{(1)} \cdot \sqrt{H}
\end{equation}

%-----------------------------------------------------------------------
\section{Stress Testing Extension}

The \texttt{stress\_means} parameter allows replacing the historical mean
vector $\boldsymbol{\mu}$ with a user-specified stress scenario mean vector
$\boldsymbol{\mu}^{\text{stress}}$, while keeping the covariance structure
unchanged:
\begin{equation}
\mathbf{r}^{(i)} \sim \mathcal{N}_N\!\left(\boldsymbol{\mu}^{\text{stress}},\, \hat{\Sigma}\right)
\end{equation}

%-----------------------------------------------------------------------
\section{Parameters}

\begin{tabular}{llll}
\toprule
Parameter & Type & Default & Description \\
\midrule
\texttt{window}  & int   & 250   & Lookback window for covariance estimation \\
\texttt{n\_sim}  & int   & 1000  & Number of Monte Carlo simulation runs \\
\texttt{distr}   & str   & \texttt{'norm'} & Distribution: \texttt{'norm'} or \texttt{'t'} \\
\texttt{decay}   & float & 0.94  & EWMA decay $\lambda$ for covariance matrix \\
\texttt{ewma\_cov} & bool & True & Use EWMA covariance (else: sample covariance) \\
\texttt{qtls}    & list  & [0.05, 0.01, 0.001] & Tail probability levels \\
\texttt{holding\_period} & int & 1 & Forecast horizon in days \\
\texttt{verbose} & bool  & True  & Print progress information \\
\bottomrule
\end{tabular}

\medskip
\noindent Additional parameters passed directly to \texttt{mc\_simulate()}:

\begin{tabular}{llll}
\toprule
Parameter & Type & Default & Description \\
\midrule
\texttt{deg\_f}      & int   & 5    & Student-$t$ degrees of freedom \\
\texttt{mlest\_cov}  & bool  & False& Use ML covariance estimator \\
\texttt{stress\_means} & list & []  & Stress scenario mean vector \\
\bottomrule
\end{tabular}

%-----------------------------------------------------------------------
\section{Advantages and Limitations}

\begin{tabular}{p{6cm} p{8cm}}
\toprule
Advantages & Limitations \\
\midrule
Captures cross-asset dependencies & Computationally intensive ($n \times T_{\text{sim}}$ draws) \\
No parametric tail assumption required & Does not capture volatility clustering \\
Flexible: Normal or Student-$t$ & Covariance matrix may be poorly conditioned \\
Stress scenario support & Convergence requires large $n$ for extreme quantiles \\
Transparent and intuitive & Square-root scaling is approximate \\
\bottomrule
\end{tabular}

%-----------------------------------------------------------------------
\section{Convergence Guidance}

For stable 99.9\% VaR estimates ($\alpha = 0.001$), at least $n = 5{,}000$
simulations are recommended. With the default $n = 1{,}000$, the 99.9\%
estimate has a coefficient of variation of approximately 10\%.

%-----------------------------------------------------------------------
\section{Usage Example}

\begin{lstlisting}[language=Python]
from risk_analytics import risk_engines as re

var_mc = re.portfolio_vaR(
    rets_orig  = rets,
    wgts       = weights,
    engine     = 'mc',
    fmt_engine = {
        'window'         : 250,
        'n_sim'          : 5000,
        'distr'          : 'norm',
        'decay'          : 0.94,
        'ewma_cov'       : True,
        'qtls'           : [0.05, 0.01, 0.001],
        'holding_period' : 1,
        'verbose'        : True,
    }
)
\end{lstlisting}

%-----------------------------------------------------------------------
\section{References}

J.P.\ Morgan/Reuters (1996). \textit{RiskMetrics Technical Document}, 4th ed.\\
Jorion, P. (2006). \textit{Value at Risk}, 3rd ed. McGraw-Hill.

\end{document}
